% Options for packages loaded elsewhere
% Options for packages loaded elsewhere
\PassOptionsToPackage{unicode}{hyperref}
\PassOptionsToPackage{hyphens}{url}
\PassOptionsToPackage{dvipsnames,svgnames,x11names}{xcolor}
%
\documentclass[
  english,
  12pt,
]{report}
\usepackage{xcolor}
\usepackage[margin=2.5cm]{geometry}
\usepackage{amsmath,amssymb}
\setcounter{secnumdepth}{5}
\usepackage{iftex}
\ifPDFTeX
  \usepackage[T1]{fontenc}
  \usepackage[utf8]{inputenc}
  \usepackage{textcomp} % provide euro and other symbols
\else % if luatex or xetex
  \usepackage{unicode-math} % this also loads fontspec
  \defaultfontfeatures{Scale=MatchLowercase}
  \defaultfontfeatures[\rmfamily]{Ligatures=TeX,Scale=1}
\fi
\usepackage{lmodern}
\ifPDFTeX\else
  % xetex/luatex font selection
\fi
% Use upquote if available, for straight quotes in verbatim environments
\IfFileExists{upquote.sty}{\usepackage{upquote}}{}
\IfFileExists{microtype.sty}{% use microtype if available
  \usepackage[]{microtype}
  \UseMicrotypeSet[protrusion]{basicmath} % disable protrusion for tt fonts
}{}
\usepackage{setspace}
% Make \paragraph and \subparagraph free-standing
\makeatletter
\ifx\paragraph\undefined\else
  \let\oldparagraph\paragraph
  \renewcommand{\paragraph}{
    \@ifstar
      \xxxParagraphStar
      \xxxParagraphNoStar
  }
  \newcommand{\xxxParagraphStar}[1]{\oldparagraph*{#1}\mbox{}}
  \newcommand{\xxxParagraphNoStar}[1]{\oldparagraph{#1}\mbox{}}
\fi
\ifx\subparagraph\undefined\else
  \let\oldsubparagraph\subparagraph
  \renewcommand{\subparagraph}{
    \@ifstar
      \xxxSubParagraphStar
      \xxxSubParagraphNoStar
  }
  \newcommand{\xxxSubParagraphStar}[1]{\oldsubparagraph*{#1}\mbox{}}
  \newcommand{\xxxSubParagraphNoStar}[1]{\oldsubparagraph{#1}\mbox{}}
\fi
\makeatother


\usepackage{longtable,booktabs,array}
\usepackage{calc} % for calculating minipage widths
% Correct order of tables after \paragraph or \subparagraph
\usepackage{etoolbox}
\makeatletter
\patchcmd\longtable{\par}{\if@noskipsec\mbox{}\fi\par}{}{}
\makeatother
% Allow footnotes in longtable head/foot
\IfFileExists{footnotehyper.sty}{\usepackage{footnotehyper}}{\usepackage{footnote}}
\makesavenoteenv{longtable}
\usepackage{graphicx}
\makeatletter
\newsavebox\pandoc@box
\newcommand*\pandocbounded[1]{% scales image to fit in text height/width
  \sbox\pandoc@box{#1}%
  \Gscale@div\@tempa{\textheight}{\dimexpr\ht\pandoc@box+\dp\pandoc@box\relax}%
  \Gscale@div\@tempb{\linewidth}{\wd\pandoc@box}%
  \ifdim\@tempb\p@<\@tempa\p@\let\@tempa\@tempb\fi% select the smaller of both
  \ifdim\@tempa\p@<\p@\scalebox{\@tempa}{\usebox\pandoc@box}%
  \else\usebox{\pandoc@box}%
  \fi%
}
% Set default figure placement to htbp
\def\fps@figure{htbp}
\makeatother



\ifLuaTeX
\usepackage[bidi=basic]{babel}
\else
\usepackage[bidi=default]{babel}
\fi
% get rid of language-specific shorthands (see #6817):
\let\LanguageShortHands\languageshorthands
\def\languageshorthands#1{}
\ifLuaTeX
  \usepackage[english]{selnolig} % disable illegal ligatures
\fi


\setlength{\emergencystretch}{3em} % prevent overfull lines

\providecommand{\tightlist}{%
  \setlength{\itemsep}{0pt}\setlength{\parskip}{0pt}}



 
\usepackage[style=apa,]{biblatex}
\addbibresource{references.bib}


\makeatletter
\@ifpackageloaded{caption}{}{\usepackage{caption}}
\AtBeginDocument{%
\ifdefined\contentsname
  \renewcommand*\contentsname{Table of contents}
\else
  \newcommand\contentsname{Table of contents}
\fi
\ifdefined\listfigurename
  \renewcommand*\listfigurename{List of Figures}
\else
  \newcommand\listfigurename{List of Figures}
\fi
\ifdefined\listtablename
  \renewcommand*\listtablename{List of Tables}
\else
  \newcommand\listtablename{List of Tables}
\fi
\ifdefined\figurename
  \renewcommand*\figurename{Figure}
\else
  \newcommand\figurename{Figure}
\fi
\ifdefined\tablename
  \renewcommand*\tablename{Table}
\else
  \newcommand\tablename{Table}
\fi
}
\@ifpackageloaded{float}{}{\usepackage{float}}
\floatstyle{ruled}
\@ifundefined{c@chapter}{\newfloat{codelisting}{h}{lop}}{\newfloat{codelisting}{h}{lop}[chapter]}
\floatname{codelisting}{Listing}
\newcommand*\listoflistings{\listof{codelisting}{List of Listings}}
\makeatother
\makeatletter
\makeatother
\makeatletter
\@ifpackageloaded{caption}{}{\usepackage{caption}}
\@ifpackageloaded{subcaption}{}{\usepackage{subcaption}}
\makeatother
\usepackage{bookmark}
\IfFileExists{xurl.sty}{\usepackage{xurl}}{} % add URL line breaks if available
\urlstyle{same}
\hypersetup{
  pdftitle={AlgaeBot: A Hybrid GraphRAG System for Domain-Specific Question Answering in Algae Research},
  pdfauthor={Filip Nový},
  pdflang={en},
  colorlinks=true,
  linkcolor={blue},
  filecolor={Maroon},
  citecolor={Blue},
  urlcolor={Blue},
  pdfcreator={LaTeX via pandoc}}


\title{AlgaeBot: A Hybrid GraphRAG System for Domain-Specific Question
Answering in Algae Research}
\author{Filip Nový}
\date{2026-04-08}
\begin{document}
\maketitle

\renewcommand*\contentsname{Table of contents}
{
\hypersetup{linkcolor=}
\setcounter{tocdepth}{2}
\tableofcontents
}
\listoffigures
\listoftables

\setstretch{1.5}
\newpage

\chapter*{Abstract}\label{abstract}
\addcontentsline{toc}{chapter}{Abstract}

\emph{{[}Write last. \textasciitilde250 words. Summarise: the problem,
your approach (hybrid GraphRAG), the evaluation framework (HOPE +
RAGAS), and your main findings.{]}}

\newpage

\chapter{Introduction}\label{introduction}

\section{Background and Motivation}\label{background-and-motivation}

\section{Research Questions}\label{research-questions}

\section{Thesis Structure}\label{thesis-structure}

\newpage

\chapter{Related Work}\label{related-work}

\section{Retrieval-Augmented
Generation}\label{retrieval-augmented-generation}

\section{Chunking Strategies}\label{chunking-strategies}

\section{Embedding Models and
Re-ranking}\label{embedding-models-and-re-ranking}

\section{Knowledge Graphs in NLP}\label{knowledge-graphs-in-nlp}

\section{Domain Context: Algae Research
Literature}\label{domain-context-algae-research-literature}

\newpage

\chapter{Data}\label{data}

\section{Corpus Description}\label{corpus-description}

\section{PDF Extraction Pipeline}\label{pdf-extraction-pipeline}

\section{Corpus Statistics}\label{corpus-statistics}

\newpage

\chapter{Methodology}\label{methodology}

\section{System Architecture
Overview}\label{system-architecture-overview}

\section{Chunking}\label{chunking}

\subsection{Recursive Text Splitting}\label{recursive-text-splitting}

\subsection{Semantic Chunking}\label{semantic-chunking}

\subsection{Recursive Semantic Chunking
(RSC)}\label{recursive-semantic-chunking-rsc}

\section{Evaluation Framework: HOPE}\label{evaluation-framework-hope}

\section{Embedding}\label{embedding}

\section{Retrieval}\label{retrieval}

\subsection{Dense Vector Retrieval}\label{dense-vector-retrieval}

\subsection{MultiQuery Expansion}\label{multiquery-expansion}

\subsection{Cross-Encoder Re-ranking}\label{cross-encoder-re-ranking}

\section{Prompt Engineering}\label{prompt-engineering}

\section{Generation}\label{generation}

\section{Knowledge Graph
Construction}\label{knowledge-graph-construction}

\subsection{Schema Design}\label{schema-design}

\subsection{Entity and Relation
Extraction}\label{entity-and-relation-extraction}

\subsection{Graph Storage and
Querying}\label{graph-storage-and-querying}

\section{Hybrid Retrieval (GraphRAG)}\label{hybrid-retrieval-graphrag}

\newpage

\chapter{Evaluation}\label{evaluation}

\section{Baseline RAG Results (RAGAS)}\label{baseline-rag-results-ragas}

\section{GraphRAG vs.~Baseline
Comparison}\label{graphrag-vs.-baseline-comparison}

\section{Ablation Analysis}\label{ablation-analysis}

\section{Qualitative Analysis}\label{qualitative-analysis}

\section{Limitations}\label{limitations}

\newpage

\chapter{Discussion}\label{discussion}

\newpage

\chapter{Conclusion}\label{conclusion}

\section{Future Work}\label{future-work}

\newpage

\chapter*{References}\label{references}
\addcontentsline{toc}{chapter}{References}

\printbibliography[heading=none]

\newpage

\chapter*{Appendix}\label{appendix}
\addcontentsline{toc}{chapter}{Appendix}

\section*{Appendix A: HOPE Metric Definitions and
Formulas}\label{appendix-a-hope-metric-definitions-and-formulas}
\addcontentsline{toc}{section}{Appendix A: HOPE Metric Definitions and
Formulas}

\section*{Appendix B: Prompt
Templates}\label{appendix-b-prompt-templates}
\addcontentsline{toc}{section}{Appendix B: Prompt Templates}

\section*{Appendix C: Chunking
Statistics}\label{appendix-c-chunking-statistics}
\addcontentsline{toc}{section}{Appendix C: Chunking Statistics}

\section*{Appendix D: System
Configuration}\label{appendix-d-system-configuration}
\addcontentsline{toc}{section}{Appendix D: System Configuration}





\end{document}
